\input{preamble.tex}
\title{A compatibility obstruction to iterating geometric HNN extensions}
\author{Human Compression project}
\date{Private draft, 11 October 2026}
\begin{document}
\maketitle

Piedade--Tranchida, arXiv:2505.24662v1, Section 5.4, proposes graphs of
coset incidence systems and asks whether the amalgamation and HNN
operations give a fundamental geometry independent of the spanning tree.
The following example isolates an earlier obstruction: locally admissible
edge maps need not remain admissible after the first HNN operation.
It concerns the stated iterative procedure, not the impossibility of
defining a different fundamental geometry.

\begin{proposition}
There is a graph with one vertex and two loops, all of whose vertex and
edge geometries are the same finite regular hypertope and all of whose
edge identifications are admissible, for which neither ordering of the
two prescribed geometric HNN operations is admissible throughout.
\end{proposition}
\begin{proof}
Let $A=C_2^3$ with basis $e_1,e_2,e_3$, and give it the standard coset
geometry with maximal parabolics
\[
 A_i=\langle e_j:j\ne i\rangle\quad(i=1,2,3).
\]
This is the join of three two-point geometries, hence is thin,
residually connected and flag-transitive, with trivial Borel.
The graph has this geometry at its sole vertex and at both loops.
On each loop take one boundary map to be the identity. Take the other
to be respectively the coordinate transpositions
$\sigma=(12)$ and $\tau=(23)$. These are actual automorphisms
permuting every indexed parabolic, so each edge map is admissible.

Perform the $\sigma$ operation first and call its stable letter $t$.
The resulting group is $G=A\rtimes_\sigma\langle t\rangle$, with
$t$ of infinite order. The type classes and maximal parabolics prescribed
by the source construction are
\[
 \{1,2\},\quad\{3\},\quad *,\qquad
 G_{\{1,2\}}=\langle e_3,t\rangle,
 \quad G_{\{3\}}=\langle e_1,e_2,t\rangle,
 \quad G_*=A.
\]
The standard parabolics contained in $A$ are exactly
\[
 A,\qquad \langle e_3\rangle,\qquad
 \langle e_1,e_2\rangle,\qquad 1.
\]
Indeed every intersection not involving $G_*$ contains the infinite-order
element $t$, and the intersections involving it are precisely the four
groups displayed. The second edge still identifies the parabolic $A$
with itself by $\tau$, but
\[
 \tau(\langle e_3\rangle)=\langle e_2\rangle
\]
is not on this list. Thus $\tau$ is not admissible for the new geometry.

If the $\tau$ operation is performed first, its two old type classes
are $\{1\}$ and $\{2,3\}$. The proper nontrivial parabolics in its
$A$ residue are $\langle e_1\rangle$ and $\langle e_2,e_3\rangle$.
The remaining map $\sigma$ sends $\langle e_1\rangle$ to
$\langle e_2\rangle$, again absent. Hence the reverse ordering also
fails. Both transpositions are their own inverses, so reversing loop
orientation does not repair either order.
\end{proof}

The abstract fundamental group still exists: it has presentation
\[
 \langle A,t,u\mid t^{-1}at=\sigma(a),\ u^{-1}au=\tau(a)
 \ (a\in A)\rangle
 \cong A\rtimes F(t,u).
\]
The obstruction is in the parabolic geometry retained after one operation.
A simultaneous construction, a common refinement of the retained incidence
data, or a global choice of type classes might avoid it; none is supplied
by this example. The source question must therefore distinguish an
unrestricted iterative recipe from a construction with additional coherent
type data. This private draft has not had an independent review and carries
no resolved-conjecture credit.
\end{document}
